\origpage{68}{ROBERTSON AND SEYMOUR}
\noindent which we describe in this section. For maximum generality we work with directed graphs, although that is not necessary for the applications. Thus each edge has a head and a tail (the same, for a loop). The underlying graph of a digraph is obtained by forgetting the directions of its edges. We shall work almost completely in terms of the underlying graph, and shall often fail to distinguish it from the digraph, for notational and other convenience. In particular, a digraph is \textit{connected} if the underlying graph is connected.

Let $G, H$ be both graphs or both digraphs. A \textit{model} $\phi$ of $H$ in $G$ assigns to each edge $e$ of $H$ an edge $\phi(e)$ of $G$, and to each vertex $v$ of $H$ a non-null connected subgraph $\phi(v)$ of $G$, such that

\hspace*{12pt}(i) the graphs $\phi(v)$ ($v \in V(H)$) are mutually vertex-disjoint, the edges $\phi(e)$ ($e \in E(H)$) are all distinct, and for $v \in V(H)$ and $e \in E(H)$, $\phi(e) \notin E(\phi(v))$

\hspace*{12pt}(ii) for $e \in E(H)$, if $e$ has head $u$ and tail $v$ (or, in the undirected case, has ends $u, v$) then $\phi(e)$ has a head in $V(\phi(u))$ and a tail in $V(\phi(v))$ (respectively, $\phi(e)$ has one end in $V(\phi(u))$ and the other in $V(\phi(v))$).

Thus, $H$ is isomorphic to a minor of $G$ if and only if there is a model of $H$ in $G$; indeed, this could be taken as the definition of when $H$ is isomorphic to a minor of $G$.

We would like to extend the idea of a minor to ``rooted'' digraphs. Intuitively, a rooted digraph is a digraph $G$, of which some vertices, $v_1, \dots, v_k$ say, are selected and called roots. If we take minors in an unrestricted sense, however, some of these roots may be deleted and some pairs of them may be identified, and this is notationally too complicated to be satisfactory. The best compromise seems to be to permit roots to become identified, but not to permit them to be deleted.

Thus, let us say that $(G, v_1, \dots, v_k)$ is a \textit{rooted digraph} if $G$ is a digraph and $v_1, \dots, v_k \in V(G)$ (not necessarily distinct). For $1 \le i \le k$, $v_i$ is the $i$ th root. Isomorphism of rooted digraphs is defined in the natural way; thus, isomorphic rooted digraphs have the same number of roots, and the isomorphism maps the $i$ th root of one to the $i$ th root of the other, for all $i$. Let $(G, v_1, \dots, v_k)$, $(H, u_1, \dots, u_k)$ be rooted digraphs. A \textit{model} of the second in the first is a model $\phi$ of $H$ in $G$ such that $v_i \in V(\phi(u_i))$ ($1 \le i \le k$), and if there is such a model, we say that $(H, u_1, \dots, u_k)$ is a minor of $(G, v_1, \dots, v_k)$. The set of all minors of $(G, v_1, \dots, v_k)$ is closed under isomorphism, and we call it the \textit{folio} of $(G, v_1, \dots, v_k)$. Clearly the folio is finite up to isomorphism. If $\delta \ge 0$ is an integer, we say that $(H, u_1, \dots, u_k)$ has \textit{detail} $\le \delta$ if $|E(H)| \le \delta$ and $|V(H) - \{u_1, \dots, u_k\}| \le \delta$. The $\delta$-\textit{folio} of $(G, v_1, \dots, v_k)$ is the set of all minors of $(G, v_1, \dots, v_k)$ with detail $\le \delta$. Its size (up to isomorphism) is bounded above by a function of $k$ and $\delta$, as is easily seen.
